%------------------------------------------------------------------------------%
\subsection{Exercícios }

\begin{enumerate}
\item Use a definição de logarítmo e as propriedades para calcular o valor de A
  \[A= \log_3{27}-\log1+ \log_{\frac{1}{5}}{125}\]
\item Assumindo que $\log2=0,3$, $\log3=0,47$ e $log5=0,69$, calcule o valor de :
\begin{enumerate}
    \item $\log(\frac{2}{5})$
    \item $\log12+\log20-\log\sqrt{5}$
    \item $\log1800-\log30$
\end{enumerate}
\item Resolva as equações logarítimicas abaixo

\begin{enumerate}
  \item $4-\log_2{(1-3x)}=3$
  \item $2\log{x}=\log{3x}+\log{(x-4)}$
  \item $2^{2x + 1} - 3.2^{x + 2} = 32$
   \item $\log_{25} (2x-1)=\frac{1}{2}$
          \item $(\log_4x)^2-2 \log_4x-3=0$
\end{enumerate}

\item Resolva as inequações abaixo
\begin{enumerate}
  \item $\ln(x^2-2x-2)\geq 0$
   \item  $\log_{\frac{1}{2}}(x^2-4) \geq \log_{\frac{1}{2}}(2x)$
   \item $\log_{\frac{1}{3}}(x^2-2) \geq \log_{\frac{1}{3}}(3x)+1$
        \item $\ln(6x+5) \geq \ln(8x-12)$
\end{enumerate}

\item Considere as funções $f(x)=\log_{2}\sqrt[3]{x} $ e  $g(x)=\frac{-x^2+9}{x+2}$

\begin{enumerate}
  \item Encontre $f\circ g$
  \item Calcule $f\circ g (-1)$
  \item Encontre o domínio da função $f\circ g$
\end{enumerate}
\item Determine o domínio de $f\circ g$, onde $f(x)=\ln(x-5)$ e $g(x)=x^3-7x-1$

\item Considere $f(x)=\ln(x-2)$

\begin{enumerate}
  \item Esboce o gráfico de $f$
  \item Justifique que $f$ é invertível
  \item Encontre a inversa de $f$ exibindo seu domínio e sua imagem
  \item Esboce o gráfico de $f^{-1}$
\end{enumerate}

    \item Faça o estudo do domínio da função $f(x)= \ln(-x^2+2x+3)$
   \item Considere a função $f(x)=\log_2(4-x^2)$ e responda o que se pede
    \begin{enumerate}
        \item Qual o domínio e a imagem de $f$?
        \item A equação $f(x)=1$ possui solução? Justifique sua resposta
        \item $f$ é injetora? Justifique sua resposta
        \item Quanto a simetria, $f$ é par, ímpar ou nem par, nem ímpar. Justifique sua resposta
        \end{enumerate}
    \item  Faça o esboço do gráfico das funções : $f(x)= \left( \frac{1}{3} \right)^x$ e $g(x)=\log_{\frac{1}{3}}(x)$,  de duas maneira diferentes A primeira levando em consideração que $g(x)=f^{-1}(x)$, a segunda marcando pontos

     \item A partir do gráfico de $y=\log_2(x)$, faça o esboço do gráfico da função $y=\log_2(x-1)+2$
  \item Sabendo que $\log_3(7x-1) = 3$ e que $\log_2(y^3 + 3) = 7$, pode-se afirmar que $\log_y(x^2 + 9)$ é igual a:
  \begin{enumerate}
    \item $6$
    \item $2$
    \item $4$
    \item $-2$
    \item $-4$
  \end{enumerate}

  \item A soma $\log_{10}\frac23+\log_{10}\frac34+\cdots+\log_{10}\frac{18}{19}+\log_{10}\frac{19}{20}$
  é igual a
  \begin{enumerate}
    \item $-\log_{10}20$
    \item $-1$
    \item $\log_{10}2$
    \item $1$
    \item $2$
  \end{enumerate}

  \item O valor da expressão $\log_3(5)\cdot\log_{125}(27)$ é:

  \begin{enumerate}
    \item $2/3$
    \item $2$
    \item $1$
    \item $3/2$
    \item um número irracional
  \end{enumerate}

  \item O conjunto solução da equação $\log_2(x + 1) + \log_2(x-3) = 5$ é:
  \begin{enumerate}
    \item $ S = \{7\}$
    \item $S = \{7, - 5\}$
    \item $S = \{17,3\}$
    \item $S = \{7/2\}$
  \end{enumerate}

  \item Na expressão $\ln 8-\ln 2 + 2\ln x = 0$, o valor de $x$ é:

  \item Considere as afirmações dadas abaixo, referentes a funções exponenciais e logarítmicas
  \begin{enumerate}
    \item[I.] $\{x\in\mathbb{R};\ln x\geq 0\}=[1,\infty)$

    \item[II.] Se $8^{-2x}=27$, então $2^{-2x}=3$

    \item[III.] $\log_{10}\sqrt{10000^3}=9/2$
  \end{enumerate}

  A alternativa correta é:
  \begin{enumerate}
    \item[(a)] Somente a afirmativa II é verdadeira
    \item[(b)] Somente a afirmativa I é verdadeira
    \item[(c)] Somente a afirmativa III é verdadeira
    \item[(d)] Somente as afirmativas I e II são verdadeiras
    \item[(e)] Somente as afirmativas I e III são verdadeiras
  \end{enumerate}


  \item Considere as afirmações dadas abaixo, referentes a funções exponenciais e logarítmicas
  \begin{enumerate}
    \item[I.] O valor de $\log_{0,25} 32$ é igual a $5/2$

    \item[II.] O valor de $x$ que satisfaz à equação $4^x-2^x = 56$ é $x = 3$

    \item[III.] As funções $y=\log (x^2)$ e $y=2\log (x)$ são iguais
  \end{enumerate}

  A alternativa correta é:
  \begin{enumerate}
    \item[(a)] Somente a afirmativa II é verdadeira
    \item[(b)] Somente a afirmativa I é verdadeira
    \item[(c)] Somente a afirmativa III é verdadeira
    \item[(d)] Nenhuma afirmativa é verdadeira
    \item[(e)] Somente as afirmativas I e III são verdadeiras
  \end{enumerate}
\item  Um vírus de um computador se espalha segundo a função \[c(t)= \frac{6500}{1+a.2^{bt}}\]
    Em que $c(t)$ é o número de computadores infectados no instante $t$ (em horas), contando a partir do momento em que a infecção foi detectada. Sabendo que $c(0)=100$, $c(3)=500$ e que $\log_2 3=1,5850$, determine as constantes $a$ e $b$ e o número de computadores infectados para $t=6$

    \item Faça o esboço do gráfico da função $f(x)=2^{x+1}$ e da sua inversa. Você deve inicialmente encontrar a inversa e depois fazer o esboço no mesmo plano cartesiano
\end{enumerate}

%------------------------------------------------------------------------------%
\subsection{Gabarito:}

\begin{enumerate}
    \item $A=0$
    \item
    \begin{enumerate}
\item $-0,39$
\item $2,015$
\item $1,76$
\end{enumerate}
\item
\begin{enumerate}
\item $x = -\frac{1}{3}$
\item $x = 6$
\item $x = 3$
\item $x = 3$
\item $x = 64$ ou $x = \frac{1}{4}$
\end{enumerate}
\item
\begin{enumerate}
\item $S=\{x \in \mathbb{R} : x \leq -1 \text{ ou } x \geq 3 \}$
\item $S=\{x \in \mathbb{R} : 2<x\leq1+\sqrt{5} \}$
\item $S=\{x \in \mathbb{R} : \sqrt{2} <x\leq2 \}$
\item $S=\{x \in \mathbb{R} : \frac{3}{2}<x\leq \frac{17}{2}\}$
\end{enumerate}
\item
\begin{enumerate}
\item $f\circ g(x) = \log_2\sqrt[3]{\frac{-x^2+9}{x+2}}$
\item $f\circ g(-1) = \log_2\sqrt[3]{\frac{8}{(-1)+2}}=\log_2\sqrt[3]{\frac{8}{1}}= \log_2 2 = 1$
\item O domínio de $f\circ g(x)$ é $x\in (-\infty,-3)\cup (-2,3)$
\end{enumerate}
\item $D =\{x \in \mathbb{R}: -2<x<-1 \,\, \text{ou} \,\,x>3 \}$

\item
\begin{enumerate}
    \item 
        \begin{center}
        \includegraphics[width=0.5\textwidth]{gabarito_aula3-7ad}
        \end{center}
    \item $f$ é apenas crescente, logo é invertível.
    \item A inversa de $f$ é $f^{-1}(x) = e^x+2$, com domínio $(-\infty,\infty)$
    \item
\end{enumerate}
\item $D =\{x \in \mathbb{R}: -1<x<3 \}$
\item
\begin{enumerate}
\item Domínio: $(-2, 2)$. Imagem: $(-\infty, \log_2 4)$
\item Sim, a equação $f(x)=1$ possui solução. As soluções são $x=-\sqrt{2}$ ou $x=\sqrt{2}$
\item Não, $f$ não é injetora, pois para todo $x$ em $(-2,2)$ temos $f(x)=f(-x)$
\item A função $f$ é par. $\log_2(4-x^2)=\log_2(4-(-x)^2)$, portanto $f(x)=f(-x)$
\end{enumerate}
\item
        \begin{center}
        \includegraphics[width=0.5\textwidth]{gabarito_aula4-10}
        \end{center}

\item
        \begin{center}
        \includegraphics[width=0.5\textwidth]{gabarito_aula4-11}
        \end{center}

\item letra b)
\item letra b)
\item letra c)
\item letra a)
\item $x =\frac{1}{2}$
\item letra e)
\item letra a)
\item $a = 64$; $b = -0,805$; $c(6) \approx2018,17$, como o número e computadores infectados é inteiro, o número de computadores infectados após $6$ horas é $2019$
\item
        \begin{center}
        \includegraphics[width=0.5\textwidth]{gabarito_aula4-20}
        \end{center}
\end{enumerate}
